\section*{Bab \ref{ch:b3:trace-element-modelling} --- Pemodelan Unsur Jejak}

\begin{solution}{exo:b3:trace-element-modelling:1}
Lherzolit spinel: $0.12\times0.40 + 0.88\times0.016 = 0.048 + 0.0141 =
0.062$. Lherzolit garnet: $0.195\times0.40 + 0.167\times0.11 +
0.639\times0.016 = 0.078 + 0.0184 + 0.0102 = 0.107$.
\end{solution}

\begin{solution}{exo:b3:trace-element-modelling:2}
$1/(0.05 + 0.02\times0.95) = 1/0.069 = \qty{14.5}{ppm}$;
$1/(0.05 + 0.10\times0.95) = 1/0.145 = \qty{6.9}{ppm}$.
\end{solution}

\begin{solution}{exo:b3:trace-element-modelling:3}
Pengayaan paling besar terjadi ketika $F \to 0$: $1/D = 100$. Dengan $D = 0$,
$C_L/C_0 = 1/F = 20$.
\end{solution}

\begin{solution}{exo:b3:trace-element-modelling:4}
$\mathrm{La/Yb} = 30/2.5 = 12.0$; $(\mathrm{La}/\mathrm{Yb})_N =
(30/0.232)/(2.5/0.163) = 129.3/15.34 = 8.4$.
\end{solution}

\begin{solution}{exo:b3:trace-element-modelling:5}
Tumpak: $C_S/C_0 = 0.01/(0.01 + 0.05\times0.99) = 0.168$. Fraksional:
$(0.95)^{99} = 0.0062$, 27 kali lebih rendah. Lelehan sesaat memiliki
$0.0062/0.01 = 0.62\,C_0$: lebih rendah daripada sumbernya.
\end{solution}

\begin{solution}{exo:b3:trace-element-modelling:6}
$D_0 = 0.58\times0.0003 + 0.27\times0.002 + 0.12\times0.053 + 0.03\times0.01
= 0.00737$; $P = 0.20\times0.0003 + 0.25\times0.002 + 0.45\times0.053 +
0.10\times0.01 = 0.0254$. Pada $F = 0.02$: non-modal $1/(0.00737 +
0.02\times0.9746) = 37.2$, modal $1/(0.00737 + 0.02\times0.9926) = 36.7$. Pada
$F = 0.10$: 9.54 dan 9.38, berselisih \qty{1.7}{\percent}.
\end{solution}

\begin{solution}{exo:b3:trace-element-modelling:7}
$D = 0.1$: tumpak $1/(0.1 + 0.09) = 5.26$; teragregasi $(1 - 0.9^{10})/0.1 =
6.51$. $D = 0.01$: tumpak $1/(0.01 + 0.099) = 9.17$; teragregasi
$(1 - 0.9^{100})/0.1 = 10.00$. Pasangan dengan $D = 0.1$ lebih berbeda: untuk
$D = 0.01$ kedua lelehan sudah memuat hampir seluruh unsur, dan keduanya cenderung ke
$1/F$.
\end{solution}

\begin{solution}{exo:b3:trace-element-modelling:8}
$D' = 0.99\times0.001 + 0.01 = 0.01099$. Lelehan pertama: $1/D' = 91$, dibandingkan
1000 pada \omterm{thm:b3:trace-element-modelling:fractional}{pelelehan fraksional}. Pada $F = 0.01$: $91\times0.99^{90.0} = 36.8$,
dibandingkan $1000\times0.99^{999} = 0.044$: lelehan pori mencegah residu
terkuras.
\end{solution}

\begin{solution}{exo:b3:trace-element-modelling:9}
Rasio dinaikkan dengan faktor $(D_{\mathrm{Sm}} + F(1 - D_{\mathrm{Sm}}))/(D_{\mathrm{La}}
+ F(1 - D_{\mathrm{La}}))$: pada $F = 0.01$, $0.0714/0.0177 = 4.0$; pada $F = 0.10$,
$0.156/0.107 = 1.46$. Rasio peka terhadap $F$ pada $F$ kecil dan hampir
tidak peka di atas sekitar \qty{10}{\percent}; juga bergantung pada $D_{\mathrm{Sm}}$,
sehingga pada mineral residu.
\end{solution}

\begin{solution}{exo:b3:trace-element-modelling:10}
$\frac1F\int_0^F \frac1D(1 - x)^{1/D - 1}\,\mathrm dx = \frac1F\bigl[-(1 -
x)^{1/D}\bigr]_0^F = [1 - (1 - F)^{1/D}]/F$, rumus dalam bab ini; jika
$(1 - F)^{1/D} \to 0$, nilainya menuju $1/F$ kali $C_0$. Dalam selisih \qty{1}{\percent}:
$(1 - F)^{100} < 0.01$, $F > 1 - 0.01^{0.01} = 0.045$, sekitar
\qty{4.5}{\percent}.
\end{solution}

\begin{solution}{exo:b3:trace-element-modelling:11}
Kemiringan $a = (10 - 6)/(20 - 10) = \qty{0.40}{ppm^{-1}}$; intersep $b = 10 -
0.40\times20 = 2.0$. $D_{\mathrm{Yb}} = 0.40\times0.44 = 0.176$;
$C_0^{\mathrm{La}} = 2.0\times0.44/0.824 = \qty{1.07}{ppm}$. $F_P = 1.07/20 =
\qty{5.3}{\percent}$, $F_Q = 1.07/10 = \qty{10.7}{\percent}$.
\end{solution}

\begin{solution}{exo:b3:trace-element-modelling:12}
$F = x_i/p_i$: klinopiroksen $0.12/0.45 = 0.27$, spinel $0.03/0.10 = 0.30$,
ortopiroksen $0.27/0.25 = 1.08$, olivin $0.58/0.20 = 2.9$; klinopiroksen
meleleh lebih dulu. Pada $F = 0.2$: $\bar D = (0.00737 - 0.0254\times0.2)/0.8 =
0.0029$. Setelah $F = 0.27$ nilai $p_i$ berubah (tidak ada lagi klinopiroksen untuk meleleh),
dan rumusnya, yang mengasumsikan $p_i$ konstan, tidak lagi berlaku.
\end{solution}

\begin{solution}{pb:b3:trace-element-modelling:1}
\textbf{1.} $D_{\mathrm{La}} = 0.12\times0.058 + 0.88\times0.001 = 0.00784$;
$D_{\mathrm{Yb}} = 0.12\times0.58 + 0.88\times0.16 = 0.210$.

\textbf{2.} $D_{\mathrm{La}} = 0.195\times0.058 + 0.167\times0.001 +
0.639\times0.001 = 0.0121$; $D_{\mathrm{Yb}} = 0.195\times0.58 +
0.167\times4.5 + 0.639\times0.16 = 0.967$. Dibulatkan: 0.0078, 0.21, 0.012, 0.97.

\textbf{3.} $D_{\mathrm{Yb}}$, dari 0.21 ke 0.97, karena garnet menahan
unsur tanah jarang berat ($D = 4.5$). Iterbium masih sangat sedikit inkompatibel,
$0.97 < 1$.

\textbf{4.} $(0.648/0.232)/(0.441/0.163) = 2.79/2.71 = 1.03$: pola datar
sekitar 2.8 kali kondrit.

\textbf{5.} $0.648/(0.0078 + 0.01\times0.9922) = \qty{36.6}{ppm}$; pada
\qty{5}{\percent}, \qty{11.3}{ppm}; pada \qty{10}{\percent}, \qty{6.05}{ppm}.

\textbf{6.} $F = (0.648/9.6 - 0.0078)/0.9922 = (0.0675 - 0.0078)/0.9922 =
0.0602$, \qty{6.0}{\percent}.

\textbf{7.} $\mathrm{Yb} = 0.441/(0.21 + 0.0602\times0.79) = \qty{1.71}{ppm}$
dan La/Yb $= 5.6$; lava A memiliki 1.7 dan 5.6: sesuai dengan sumber spinel.

\textbf{8.} La $= 0.0078\times9.6 = \qty{0.075}{ppm}$, Yb $= 0.21\times1.71 =
\qty{0.36}{ppm}$. Lantanum tersisa: $0.940\times0.075/0.648 = 0.11$, sekitar
\qty{11}{\percent}.

\textbf{9.} $F = 0.648/9.6 = \qty{6.75}{\percent}$.

\textbf{10.} $(0.648/0.0078)(1 - F)^{127.2}$: \qty{73}{ppm} pada
\qty{0.1}{\percent}, \qty{23}{ppm} pada \qty{1}{\percent}, \qty{0.03}{ppm} pada
\qty{6}{\percent}.

\textbf{11.} $0.648\times(1 - 0.940^{128.2})/0.060 = \qty{10.8}{ppm}$.

\textbf{12.} $F = \qty{6.75}{\percent}$: pada $F$ ini, $(1 - F)^{1/D} =
\num{1.3e-4}$, sehingga lelehan gabungan menampung seluruh lantanum sumbernya, dan
$\bar C_L = C_0/F$ seperti untuk unsur yang sepenuhnya inkompatibel.

\textbf{13.} $0.648\times0.9325^{127.2} = \qty{9e-5}{ppm}$, sekitar 800 kali
lebih sedikit daripada \qty{0.075}{ppm} pada \omterm{thm:b3:trace-element-modelling:batch}{pelelehan tumpak}.

\textbf{14.} Kedua model itu membatasi jawabannya: dari \qty{6.0}{\percent} (tumpak)
sampai \qty{6.75}{\percent} (fraksional gabungan). Lantanum menetapkan $F$ hingga
sekitar \qty{1}{\percent} dari pelelehan, apa pun modelnya.

\textbf{15.} A: La/Yb $= 5.6$, $(\mathrm{La}/\mathrm{Yb})_N =
(9.6/0.232)/(1.7/0.163) = 4.0$. B: La/Yb $= 34.7$, $(\mathrm{La}/\mathrm{Yb})_N
= (15.6/0.232)/(0.45/0.163) = 24.4$.

\textbf{16.} Ketika $F \to 0$: $(0.648/0.0078)/(0.441/0.21) = 83.1/2.10 = 39.6$.

\textbf{17.} $F = \qty{0.11}{\percent}$, saat lelehan akan memuat
\qty{72}{ppm} La dan \qty{2.1}{ppm} Yb. Lava B memiliki \qty{15.6}{ppm} La: sumber
spinel tidak mampu menghasilkannya.

\textbf{18.} $F = (0.648/15.6 - 0.012)/0.988 = 0.0299$, \qty{3.0}{\percent};
$\mathrm{Yb} = 0.441/(0.97 + 0.0299\times0.03) = \qty{0.454}{ppm}$, sesuai
dengan pengamatan.

\textbf{19.} Garnet menggantikan spinel pada sekitar \qty{60}{km} sampai \qty{90}{km}: B mulai
meleleh lebih dalam dari itu, A lebih dangkal, di dalam medan spinel. Lelehan
garnet-lherzolit dengan \qty{9.6}{ppm} La ($F \approx \qty{5.6}{\percent}$) akan memiliki La/Yb
sekitar 21, bukan 5.6: A tidak berasal dari garnet lherzolit.

\textbf{20.} $a = 0.97/0.441 = \qty{2.20}{ppm^{-1}}$, $b = 0.648\times0.03/0.441
= 0.044$; pada La $= 15.6$, La/Yb $= 2.20\times15.6 + 0.044 = 34.4$, dibandingkan
34.7 untuk lava B.

\textbf{21.} Olivin menolak kedua unsur itu: La/Yb tidak berubah. Lelehan primer
memiliki $9.6\times0.9 = \qty{8.64}{ppm}$ La, sehingga $F = (0.648/8.64 -
0.0078)/0.9922 = \qty{6.8}{\percent}$, bukan 6.0.

\textbf{22.} Satu sumber dapat menghasilkan kedua lava: B adalah lelehan kecil
(\qty{3}{\percent}) yang mulai di medan garnet, lebih dalam daripada sekitar \qty{60}{km} sampai
\qty{90}{km}; A lelehan yang lebih besar di medan spinel.

\textbf{23.} Lava A adalah \textbf{\qty[text-series-to-math=true]{6.0}{\percent}}
lelehan tumpak dari spinel lherzolit (\qty{6.75}{\percent} jika lelehan diekstraksi secara fraksional dan digabungkan).
tersaring secara fraksional dan digabungkan).
\end{solution}
